\documentclass{article}
\usepackage[utf8]{inputenc}

\usepackage{amsmath, amssymb, amsthm, amsfonts}
\usepackage{subcaption, enumerate}
\usepackage{bbold}
\usepackage{polynom}
\usepackage{xfrac}
\usepackage{graphicx}
\usepackage{mathtools}
\usepackage{enumitem}

\usepackage{tikz}

\newtheorem{theorem}{Theorem}
\newtheorem{corot}{Corollary}[theorem]
\newtheorem{lemma}{Lemma}
\newtheorem{corol}{Corollary}[lemma]
\newtheorem{proposition}{Proposition}
\newtheorem{corop}{Corollary}[proposition]

\theoremstyle{definition}
\newtheorem{definition}{Definition}
\newtheorem{exampled}{Example}[definition]
\newtheorem{examples}{Example}[section]
\newtheorem*{notation}{Notation}

\newcommand{\dab}
{
	\begin{tikzpicture}[scale=0.04]
		% head
		\draw[thick] (0, 0) circle (3cm);

		% bent arm
		\draw[thick, rotate around={15 : (0, 0)}] (-3, 0) -- (-3.7, -2) -- (4.1, -1.1);

		% extended arm
		\draw[thick, rotate=20] (3, 0) -- (7, 2);
	\end{tikzpicture}
}
\newcommand{\st}{\text{ s.t. }}

\makeatletter
\renewcommand*\env@matrix[1][*\c@MaxMatrixCols c]{%
  \hskip -\arraycolsep
  \let\@ifnextchar\new@ifnextchar
  \array{#1}}
\makeatother

\usepackage{fullpage}

\title{Topology \\ with Dr Nevard}
\author{Laura Gallo}
\date{November 12, 2020}

\begin{document}
\maketitle

\section{Topologies and Metrics}
\begin{definition}
	Let $X$ be a set. A collection of subsets of $x$, $\mathcal{T}$, is a \emph{topology} on $X$ if:
	\begin{enumerate}
		\item $\emptyset \in \mathcal{T}$
		\item if $U_1, U_2 \in \mathcal{T}, U_1 \cap U_2 \in \mathcal{T}$
		\item if $\{U_\alpha\} \subset \mathcal{T}, \bigcup_\alpha U_\alpha \in \mathcal{T}$
	\end{enumerate}
	\noindent
	Elements of $\mathcal{T}$ are called \emph{open sets}.
\end{definition}

\begin{exampled}
	For $\mathbb{R}$, the set of all open intervals $(a, b)$ and their unions is a topology. This is the standard topology on $\mathbb{R}$
\end{exampled}
\begin{exampled}
	For $\mathbb{R}^n$, the set of all open intervals $(a_1, b_1) \times (a_2, b_2) \times \dots \times (a_n, b_n)$ and their unions is a topology. This is the standard topology on $\mathbb{R}^n$
\end{exampled}

\begin{definition}
	A metric on a topology $\mathcal{T}$ is a function $\mu: \mathcal{T}\times\mathcal{T} \rightarrow [0, \infty)$ such that
	\begin{enumerate}
		\item $\mu(x, y) = 0 \Leftrightarrow x=y$
		\item $\mu(x, y) = \mu(y, x)$
		\item $\mu(x, y) \leq \mu(x, z) + \mu(z, y)$
	\end{enumerate}
\end{definition}

\begin{exampled}
	$\mathbb{R}$, $d(x, y)=|x-y|$
\end{exampled}
\begin{exampled}
	$\mathbb{R}^n$, $d(x, y)=\sqrt{\sum_{i=0}^n (x_i-y_i)^2}$
\end{exampled}
\begin{exampled}
	$\ell^\infty$ (the set of infinite sequences), $\sup_{i\in \mathbb{N}}|x_i-y_i|$
\end{exampled}
\begin{exampled}
	$L^\infty$, $\sup_{x\in\text{ domain}} |f(x)-g(x)|$
\end{exampled}
\begin{exampled}
	$L^p$, $(\int_X |f(t)-g(t)|^pdt)^{\frac{1}{p}}$
\end{exampled}
	\noindent
	The above examples have metrics that are also norms on the given vector space; however, a metric does not necessarily have to be a norm. The below examples are valid metrics on the standard topology on $\mathbb{R}$:
\begin{exampled}
	$d(x, y) = \frac{|x-y|}{1 + |x-y|}$
\end{exampled}
\begin{exampled}
	\begin{math}
		d(x, y) =
		\begin{cases}
			1 & \text{if $x \neq y$} \\
			0 & \text{if $x = y$}
		\end{cases}
	\end{math}
\end{exampled}

\section{Neighbourhoods}
If we have a collection of open sets $\{U_\alpha\}$, $x\in U_\alpha$ means $U_\alpha$ is a neighbourhood of $x$. If $U_\alpha \subset U_\beta$, $U_\alpha$ is a smaller neighbourhood. We can say that $y_n \rightarrow x$ if $\forall$ neighbourhoods $U_\alpha$ of $x$, given $U_{\alpha_0}$, $\exists N$ s.t. $n > N \Rightarrow y_n \in U_\alpha$.

\section{Bases}
\begin{definition}
	A base for a topology $\mathcal{T}$ is a subset $\mathcal{B} \subset \mathcal{T}$ such that the set of all unions of $\mathcal{B}$ is $\mathcal{T}$.
\end{definition}

\begin{exampled}
	The standard topology on $\mathbb{R}$ has the countable base $\mathcal{B}=\{(r, s) : r, s \in \mathbb{Q}\}$
\end{exampled}
\begin{exampled}
	The standard topology on $\mathbb{R}^n$ has a similar countable base $\mathcal{B}=\{(a_1, b_1)\times(a_2, b_2)\times\dots\times(a_n, b_n)\}$.
\end{exampled}
\begin{exampled}
	The standard topology on $\mathbb{R}^n$ has another base $\mathcal{B}_r(x)=\{y\in\mathbb{R}^n:||x-y||<r\}$
\end{exampled}

\noindent
If $X$ is a normed space with norm $||\cdot||$, then the base for the topology of $X$ is $\mathcal{B}_r(x)=\{y\in X:||x-y||<r\}$.

\section{Equivalent Metrics/Norms}
\begin{definition}
	Two metrics/norms are equivalent if they yield identical topologies.
\end{definition}
\begin{exampled}
	In $\mathbb{R}$, $d_1(x, y)=|x-y|$ and $d_2(x, y)=\frac{|x-y|}{1+|x-y|}$ are equivalent. This is left as an exercise to the reader.
\end{exampled}%%%%

\end{document}


